\chapter*{Phân Công Công Việc}
\addcontentsline{toc}{chapter}{Phân Công Công Việc}

Nhóm chia đồ án thành ba mảng có khối lượng tương đương, mỗi mảng gồm cả phần
thiết kế thuật toán, phần cài đặt và phần kiểm thử tương ứng. Mỗi thành viên là
người chịu trách nhiệm chính cho mảng của mình: viết mã, viết test, đo số liệu và
soạn các mục báo cáo liên quan. Việc rà soát chéo mã nguồn được thực hiện cho
toàn bộ ba mảng.

\footnotesize
\begin{longtable}{>{\raggedright\arraybackslash}p{2.8cm}>{\raggedright\arraybackslash}p{5.75cm}p{6.15cm}}
\caption*{\textbf{Bảng phân công nhiệm vụ giữa các thành viên}}\\
\toprule
\textbf{Thành viên} & \textbf{Nội dung phụ trách chính} & \textbf{Sản phẩm bàn giao} \\
\midrule
\endfirsthead
\toprule
\textbf{Thành viên} & \textbf{Nội dung phụ trách chính} & \textbf{Sản phẩm bàn giao} \\
\midrule
\endhead
Hoàng Tiến Đạt\newline 23010864
&
\textbf{Tầng tối ưu lộ trình.} Cài đặt Dijkstra nhiều đích trên MultiDiGraph,
xây ma trận chi phí có hướng, quy hoạch động Held--Karp cho ba kiểu kết thúc,
hai heuristic nearest-neighbor và 2-opt, thuật toán khôi phục tuyến chi tiết.
Thiết kế và chạy benchmark tầng thứ tự ghé.
&
\texttt{src/shortest\_path.py},\newline
\texttt{src/tsp.py},\newline
\texttt{src/heuristics.py},\newline
\texttt{src/route\_optimizer.py},\newline
\texttt{scripts/benchmark\_tsp.py},\newline
\texttt{tests/test\_shortest\_path.py},\newline
\texttt{tests/test\_tsp.py},\newline
\texttt{tests/test\_heuristics.py};\newline
Mục 2.2--2.3, 4.3--4.4, 5.4 \\
\addlinespace
Phạm Văn Tuấn Kiệt\newline 23010881
&
\textbf{Mô hình truy tìm thông tin ẩn.} Thiết kế luật trò chơi và thứ tự ưu tiên
kết thúc, kiểu \texttt{Observation} tách ranh giới thông tin, cơ chế lan truyền
tập niềm tin, ba chiến thuật truy tìm và hàm điểm phối hợp. Thiết kế và chạy
benchmark 360 ván.
&
\texttt{src/chase.py},\newline
\texttt{src/chase\_maps.py},\newline
\texttt{scripts/benchmark\_chase.py},\newline
\texttt{scripts/prepare\_chase\_map.py},\newline
\texttt{tests/test\_chase.py},\newline
\texttt{tests/test\_chase\_maps.py};\newline
Mục 2.4--2.6, 3.5, 4.5--4.6, 5.5--5.7 \\
\addlinespace
Mầu Danh Chiến\newline 23017270
&
\textbf{Dữ liệu, giao diện và chất lượng.} Chuẩn bị và đóng gói graph Hà Nội,
quản lý tệp lớn bằng Git LFS, lớp nạp/cache và kiểm tra metadata, giao diện
Streamlit hai chế độ, hiển thị bản đồ Folium, replay và xuất log JSON. Kiểm thử
tích hợp, kiểm tra trình duyệt thật, biên soạn báo cáo \LaTeX{} và sinh hình vẽ.
&
\texttt{app.py},\newline
\texttt{src/map\_service.py},\newline
\texttt{src/chase\_ui.py},\newline
\texttt{src/visualization.py},\newline
\texttt{scripts/preload\_hanoi*.py},\newline
\texttt{scripts/make\_report\_figures.py},\newline
\texttt{tests/test\_app*.py},\newline
\texttt{tests/test\_services.py};\newline
Chương 1, 3, 6, 7 và phụ lục \\
\addlinespace
\midrule
Cả nhóm
&
Thống nhất mô hình toán và giao diện giữa các mô-đun, rà soát chéo mã nguồn,
chạy lại toàn bộ 112 ca kiểm thử trước mỗi mốc bàn giao, dựng kịch bản demo và
chuẩn bị nội dung bảo vệ.
&
Tóm tắt, Abstract, kịch bản demo,\newline bộ số liệu trong \texttt{output/} \\
\bottomrule
\end{longtable}
\normalsize

Khối lượng được cân bằng theo ba tiêu chí: số mô-đun phụ trách, số ca kiểm thử
phải viết và số mục báo cáo phải soạn. Hai mảng thuật toán (Đạt và Kiệt) nặng về
chứng minh tính đúng và phân tích độ phức tạp; mảng còn lại (Chiến) nặng về xử lý
dữ liệu thật, tích hợp và bảo đảm chất lượng. Ba mảng đều có phần đo đạc riêng và
đều cần thiết để sản phẩm chạy được đầu--cuối.
\clearpage
